\Article{23}
Whenever experiments are carried out to determine the absolute value $T$ of terrestrial magnetism, it is important to complete them all within a fairly short time, so that there is no reason to fear any appreciable change in the magnetic state of the needles used. In observing the deflections of the movable needle, it is therefore advisable to follow only the first method of Article $20$, using just \SourceLeaf{116}two distinct, suitably chosen distances, provided two terms of the series suffice. Of the several applications of this method, we choose as an example the one to which we devoted the most scrupulous care, measuring the distances with microscopic precision.

The experiments were carried out on $1832$ Sept. $18$ with two instruments, which we will distinguish by the letters $A,B$, and with three needles, which we will distinguish by the numbers $1,2,3$. Needles $1,2$ are the ones called the first and second in Article $11$. The experiments fall into two parts.

First the simultaneous oscillations of needle $1$ in instrument $A$ and needle $2$ in instrument $B$ were observed. The time for one oscillation, reduced to infinitesimal amplitude, was
\begin{center}
\begin{tabular}{lcl}
needle $1$&.\ .\ .\ .\ .\ .&$15''22450$\\
needle $2$&.\ .\ .\ .\ .\ .&$17{,}29995$
\end{tabular}
\end{center}
The former value was inferred from $305$ oscillations and the latter from $264$.

Next needle $3$ was suspended in instrument $A$, while needle $1$ was placed on a straight line perpendicular to the magnetic meridian, both to the east and to the west, in two ways on each side. The deflection of needle $3$ was observed for each position of needle $1$. Repeating these experiments at two different distances $R$ yielded the following values of the angle $v$, understood as in Articles $19$, $21$.
\[
\begin{aligned}
R&=1^{\mathrm m}2,&\qquad v&=3^0\ 42'\ 19''4\\
R'&=1{,}6.&v'&=1\ 34\ 19{,}3
\end{aligned}
\]
The oscillations of needle $2$ in instrument $B$ were also observed during these experiments. The time for one infinitesimal oscillation corresponding to the mean time was inferred from $414$ oscillations and was $=17''29484$.

The times were measured by a chronometer whose daily loss was $=14''24$.

Let $M,m$ denote the moments of free magnetism for needles $1$ and $3$, and let $\theta$ be the torsion constant of the thread in instrument $A$ while it supported needle $1$ or $3$ (their weights are nearly the same). We have
\[
\frac{\theta}{TM}=\frac1{597{,}4}
\]
as in Article $11$, and
\[
\frac{\theta}{Tm}=\frac1{721{,}6}
\]
since needle $3$ was more strongly magnetized than needle $1$.
\par
\SourceLeaf{117}
The moment of inertia of needle $1$ was already known from earlier experiments (see Article $11$). They had yielded $K=4228732400$ with the millimeter and the milligram as units.

The thermometer reading in each of the two rooms where the instruments were set up varied so little throughout the experiments that there is no need to take it into account.

Let us now calculate the results of these experiments to obtain the intensity $T$ of terrestrial magnetism. The variation in the oscillations of needle $2$ reveals a slight change in this intensity. To arrive at a definite value, we will therefore reduce the observed oscillation time of needle $1$ to the mean state of terrestrial magnetism during the second part of the observations. This time requires a further correction for the chronometer's loss and a third for the torsion of the thread. This gives the reduced time of one oscillation of needle $1$
\[
\begin{aligned}
&=15{,}22450\times\frac{17{,}29484}{17{,}29995}\mathpunct{.}\frac{86400}{86385{,}76}\mathpunct{.}\BareRoot{\frac{598{,}4}{597{,}4}}\\
&=15''23530=t
\end{aligned}
\]

From this we obtain the product $TM=\frac{\pi\pi K}{tt}=179770600$. The slight difference between this value and the one we found above in Article $11$ for September $11$ is due both to the change in terrestrial magnetism and to the change in the magnetic state of the needle.

From the observed deflections we obtain
\[
F=\frac{R'^5\operatorname{tang}v'-R^5\operatorname{tang}v}{R'R'-RR}=113056200
\]
if we take the millimeter as the unit. From this we obtain
\[
\frac MT=\frac12F\left(1+\frac{\theta}{Tm}\right)=56606437
\]
Comparing this number with the value of $TM$ finally gives
\[
T=1{,}782088
\]
for the intensity of the horizontal terrestrial magnetic force on September $18$ at $5$ o'clock.

